\chapter{Método dos Elementos Finitos}

% Configurações de documento e pacotes equivalentes

\begin{figure}[htbp]
  \centering
  \begin{tikzpicture}[
      node distance=1.5cm and 1.0cm,
      block/.style={rectangle, draw=blue!70, fill=blue!5, rounded corners=5pt, thick, align=center, inner sep=6pt},
      line/.style={draw, -latex, thick, color=gray!80}
    ]
    \node [block, text width=6cm] (step1) {\textbf{1. Definição do Problema} \\ \small Vaso de pressão multicamadas};
    \node [block, text width=5.5cm, below=0.8cm of step1] (step2) {\textbf{2. Parâmetros Iniciais} \\ \small Geometria, Fração volumétrica, Condições de contorno};
    \node [block, text width=4.8cm, below left=1.0cm and 0.2cm of step2] (step3a) {\textbf{3a. Modelo Analítico} \\ \small Homogeneização Assintótica (MHA)};
    \node [block, text width=4.8cm, below right=1.0cm and 0.2cm  of step2] (step3b) {\textbf{3b. Modelo Numérico} \\ \small Elementos Finitos (MEF -- FEniCS)};
    \node [block, text width=5.5cm, below=3.2cm of step2] (step4) {\textbf{4. Verificação} \\ \small Comparação (Tensões e Deslocamentos)};
    \node [block, text width=6cm, below=0.8cm of step4] (step5) {\textbf{5. Análise de Viabilidade} \\ \small Resistência ($W_g$) e Rigidez ($K_g$)};

    \draw [line] (step1) -- (step2);
    \draw [line] (step2) -| (step3a);
    \draw [line] (step2) -| (step3b);
    \draw [line] (step3a) |- (step4);
    \draw [line] (step3b) |- (step4);
    \draw [line] (step4) -- (step5);
  \end{tikzpicture}
  \caption{Fluxograma geral da metodologia adotada. Fonte: Autor.}
  \label{fig-fluxograma}
\end{figure}

A modelagem numérica de estruturas elásticas microperiódicas com alto contraste de propriedades exige uma transição rigorosa entre as formulações matemáticas contínuas desenvolvidas no Capítulo 3 e a sua representação no domínio discreto. Este capítulo expõe a metodologia prática e o fluxo computacional adotados para a simulação de vasos de pressão cilíndricos multicamadas por meio do Método dos Elementos Finitos (MEF). São detalhados a estratégia de discretização espacial do domínio axissimétrico, a lógica de programação para a geração da microestrutura bimetálica periódica via FEniCS, o equacionamento das condições de contorno de pressão e simetria longitudinal, e o protocolo de cálculo de métricas de eficiência estrutural baseados no código computacional real de projeto implementado no arquivo principal \texttt{main.py}.

\section{Discretização e Configuração da Malha}\label{sec:discretizacao_malha}

A precisão de um solucionador de elementos finitos para problemas com heterogeneidades de alta frequência e gradientes de tensão severos é governada pelas características geométricas do domínio computacional e pelo refinamento espacial da malha de discretização \cite{Sadd2009}. Esta seção detalha a definição do domínio, a escolha da classe de elementos finitos e o equacionamento para o posicionamento exato dos nós nas interfaces de transição elástica.

\subsection{Geometria do Domínio Computacional Axissimétrico}

Conforme discutido na formulação teórica da elasticidade de revolução (Seção 3.4), a invariância circunferencial ($\partial_{\theta} = 0$) e a hipótese de deformação plana generalizada reduzem o problema elástico tridimensional (3D) de um cilindro oco para uma seção transversal bidimensional (2D) contida no plano radial-axial $r-z$ \cite{Timoshenko_1951}.

Para manter total consistência físico-matemática e dimensional com o código numérico desenvolvido, o domínio computacional espacial, denotado por $\Omega_{rz}$, é representado geometricamente por uma seção retangular (ver figura ...)  plana cujas dimensões de projeto em escala real de simulação são dadas por:

%TODO: Adicionar figura da seção retangular, mostrando eixos r e z

\begin{equation}
  \Omega_{rz} = \{ (r, z) \in \mathbb{R}^2 : R_{\text{int}} \le r \le R_{\text{ext}} , \quad 0 \le z \le H \}
  \label{eq:dominio_mef}
\end{equation}

onde:

\begin{itemize}
  \item $R_{\text{int}} = 9,0 \, \text{mm}$ representa o raio interno do cilindro de referência;
  \item $R_{\text{ext}} = 11,0 \, \text{mm}$ representa o raio externo do cilindro;
  \item A largura do domínio equivale à espessura total da parede do vaso, dada por $t = R_{\text{ext}} - R_{\text{int}} = 2,0 \, \text{mm}$;
  \item A dimensão axial (altura do retângulo de revolução) é fixada em $H = 2,0 \, \text{mm}$.
\end{itemize}

%Essa escala geométrica em milímetros é rigorosamente compatível com o sistema de unidades (N, mm, MPa) adotado no script principal, no qual as constantes constitutivas elásticas são prescritas em megapascals e a pressão nominal aplicada ($30 \, \text{MPa}$) coincide numericamente com $30 \, \text{N/mm}^2$, evitando inconsistências de fator de escala numérico durante a montagem das equações elásticas.

% REMOVE: Não acho esse trecho necessário, pois essa escolha de unidades é bastante conhecida.

\subsection{Escolha do Elemento Finito (Lagrange de Grau 1)}

Na implementação do solucionador no FEniCS, as componentes do vetor de deslocamento $\mathbf{u} = (u_r, u_z)$ e o campo de funções de teste virtuais $\mathbf{v} = (v_r, v_z)$ são interpolados espacialmente utilizando \textbf{elementos finitos de Lagrange de grau 1} ($P_1$) também chamados de \textit{Constant strain elements} (CST) \cite{Sadd2009}, configurando um espaço de funções contínuas e lineares de Galerkin \cite{logg2012automated}.

A adoção de polinômios lineares, em conjunto com uma malha altamente refinada e estruturada, é uma premissa metodológica fundamental para a simulação elástica eficiente de cilindros concêntricos sob pressão pelas seguintes razões elásticas e computacionais:

% 1. \textbf{Invariância de Deformação e Tensão por Elemento:} Como a deformação $\mathbf{\epsilon}$ depende das derivadas primeiras do campo de deslocamento ($\mathbf{\epsilon} = \frac{1}{2}(\nabla \mathbf{u} + \nabla^T \mathbf{u})$), o uso de funções de interpolação lineares ($P_1$) garante que as deformações e as tensões no interior de cada elemento finito individual sejam constantes ao longo de cada triângulo \cite{Sadd2009}. Isso é perfeitamente congruente com a física de meios laminados periódicos (onde o material e as propriedades elásticas também são assumidos como constantes por trecho/elemento), fornecendo degraus de tensões constantes por subcamada elástica sem introduzir suavizações ou gradientes parabólicos espúrios na interface de descontinuidade elástica.

% REMOVE: Esse trecho não está claro para mim. Não acho que fato de que em elementos P1 os campos de deformação e tensão serem constantes necessariamente ajude na simulação só pelo fato das propriedades serem constantes.

1. \textbf{Minimização Drástica do Custo Computacional:} Como o presente estudo conduz simulações com um elevado número de células periódicas (até $N = 200$) e resoluções de malha extremamente refinadas, a adoção de elementos $P_1$ limita significativamente o número total de graus de liberdade globais a serem resolvidos. Isso reduz drasticamente o tempo de CPU e os requisitos de memória RAM, permitindo a execução de varreduras paramétricas e tabelas de erros quase que instantaneamente.
2. \textbf{Resolução Radial por Refinamento Dinâmico:} A potencial limitação de menor precisão local de elementos lineares em gradientes acentuados é completamente mitigada pelo emprego de uma altíssima densidade radial de elementos finitos, garantindo que as curvas globais de deslocamento e tensão sejam perfeitamente discretizadas de forma suave.

\subsection{Configuração de Alta Discretização e Densidade Radial}

A malha do domínio $\Omega_{rz}$ é estruturada de forma mapeada, contendo $N_z$ divisões na direção axial e $N_r$ divisões na direção radial, resultando em uma malha de elementos triangulares obtidos pela divisão diagonal do retângulo espacial \cite{logg2012automated}. Essa discretização ao longo do raio $r$ é governada pela resolução espacial requerida para capturar os degraus microscópicos locais de descontinuidade constitutiva nas subcamadas.

% TODO: Adicionar figura mostrando a malha no domínio 2D, com as camadas de material e as fronteiras de referência. (Não esquecer de usar as cores)

No algoritmo implementado no arquivo principal \texttt{main.py}, a densidade de elementos radiais respeita de forma exata e proporcional as frações volumétricas elásticas de cada camada, definindo-se no script:

\begin{itemize}
  \item \texttt{V1}: Fração volumétrica do Material 1;
  \item \texttt{V2}: Fração volumétrica do Material 2;
  \item \texttt{N\_cel}: Número de células periódicas;
  \item \texttt{resolucao\_malha}: Número de elementos radiais por célula periódica;
  \item \texttt{nel1 = int(V1 * resolucao\_malha)}: número de elementos na camada do Material 1 (Aço) por célula periódica;
  \item \texttt{nel2 = int(V2 * resolucao\_malha)}: número de elementos na camada do Material 2 (Alumínio) por célula periódica, com $V_2 = 1 - V_1$;
  \item \texttt{nel\_celula = nel1 + nel2}: número total de elementos radiais contidos em cada período elástico básico (célula unitária).
\end{itemize}

Dessa forma, o número total de elementos na direção radial ($N_r$, ou \texttt{n\_elements\_r}) e axial ($N_z$, ou \texttt{n\_elem\_z}) são estabelecidos diretamente no gerador \texttt{RectangleMesh}:

\begin{equation}
  N_r = N_{\text{cel}} \times \text{resolucao\_malha}
  \label{eq:nr_mef_v3}
\end{equation}

\begin{equation}
  N_z = n_{\text{elem\_z}}
  \label{eq:nz_mef_v3}
\end{equation}

\begin{verbatim}[language=Python, caption={Trecho de discretização espacial do domínio de parede do vaso (Implementado no script principal).}]
# Trecho de discretização espacial do domínio de parede do vaso
nel1 = int(V1 * resolucao_malha)
nel2 = int(V2 * resolucao_malha)
nel_celula = nel1 + nel2
n_elements_r = N_celulas * nel_celula

mesh = RectangleMesh(
    Point(R_int, 0.0), Point(R_ext, H),
    int(n_elements_r), int(n_elem_z)
) 
\end{verbatim}

% TODO: Verificar se nel_celula pode ser removido pois é igual a resolucao_malha.s

Adotando-se a configuração padrão de simulação com $N_{\text{cel}} = 200$ células periódicas, \texttt{resolucao\_malha} = 20 e $N_z = 10$, a malha radial totaliza \textbf{$N_r = 4.000$ elementos radiais}, gerando uma discretização de $80.000$ elementos triangulares no domínio bidimensional. Essa altíssima densidade radial garante que cada subcamada individual de aço ou alumínio (cuja espessura de período para $N=200$ é de apenas $5,0 \, \mu\text{m}$ no aço e $5,0 \, \mu\text{m}$ no alumínio) contenha exatamente \textbf{10 elementos lineares internos}, permitindo mapear com riqueza de detalhes elásticos as flutuações e tensões microscópicas.

\subsection{O Alinhamento dos Nós com as Interfaces Físicas}

%No Método dos Elementos Finitos, a distribuição das propriedades mecânicas do material (módulo de elasticidade $E$ e coeficiente de Poisson $\nu$) é avaliada e integrada numericamente nos pontos de integração Gaussiana localizados no interior de cada elemento finito \cite{logg2012automated}.

%Se um elemento finito cruzar a interface que divide uma camada de aço de uma camada de alumínio, as propriedades elásticas integradas para aquele elemento sofrerão uma média suave artificial ("borramento" interfacial). Como consequência física direta, o salto descontínuo na tensão circunferencial ($\sigma_{\theta\theta}$) e a mudança de inclinação (\textbf{kinks}) no deslocamento radial ($u_r$), característicos da interface de contato perfeito bimetálica, seriam suavizados, arruinando completamente a acurácia física e a consistência local da simulação numérica.  

No Método dos Elementos Finitos, o deslocamento é encontrado nos nós e interpolado para o interior de cada elemento finito \cite{logg2012automated}.

Se um elemento finito cruzar a interface que divide uma camada de aço de uma camada de alumínio, o deslocamento será interpolado de forma suave entre as duas camadas, o que pode levar a resultados imprecisos.

Para blindar o modelo contra essa perda de precisão, a fronteira de cada subcamada material \textbf{deve coincidir estritamente com as bordas dos elementos finitos}. Ou seja, os nós da malha espacial devem estar localizados exatamente sobre as coordenadas radiais das interfaces elásticas.

Essa restrição é perfeitamente satisfeita no código real ao impor o arredondamento inteiro de \texttt{nel1} e \texttt{nel2} por célula periódica através da função \texttt{int()} em Python. Ao garantir que $n_{\text{el\_cam1}}, n_{\text{el\_cam2}} \in \mathbb{N}$, a espessura de cada camada física se torna um múltiplo inteiro exato do tamanho local do elemento ($h_r = t / N_r$). O gerador de malhas realiza a transição de propriedades de forma abrupta a cada nó interfacial, eliminando oscilações numéricas próximas às interfaces e garantindo a convergência estrita da simulação heterogênea para o limite assintótico homogeneizado do MHA.

A simulação numérica de meios altamente heterogêneos e microperiódicos, como o cilindro bimetálico composto por camadas alternadas de aço carbono e liga de alumínio, impõe desafios significativos de modelagem geométrica. À medida que o número de períodos $N_{\text{cell}}$ aumenta, a espessura de cada subcamada individual encolhe de forma inversamente proporcional, tornando a reconstrução física de cada interface no ambiente de CAD tradicional e o posterior corte da malha (\textit{mesh partitioning}) processos manuais altamente ineficientes e propensos a erros de precisão numérica. Neste trabalho, optou-se por não automatizar a geração de uma malha conforme às interfaces físicas; em vez disso, adota-se o mapeamento implícito de propriedades descrito a seguir.

Para contornar essa barreira operacional, adota-se neste trabalho a estratégia de \textbf{Mapeamento Implícito de Propriedades} via vetor de interfaces discretas. Em vez de discretizar as interfaces físicas na própria malha, o domínio espacial é particionado por uma malha estruturada única e contínua no plano $r-z$ (conforme já discutido na Seção~\ref{sec:discretizacao_malha}). As propriedades elásticas locais dos materiais de fase (módulo de Young, coeficiente de Poisson e a respectiva matriz constitutiva) são injetadas diretamente nos pontos de integração Gaussiana através de expressões analíticas espaciais periódicas de busca indexada programadas no ambiente FEniCS \cite{logg2012automated}.

% TODO: Verificar o que é "Mapeamento Implícito de Propriedades" e adicionar uma citação, se possível.

\subsection{Formulação Matemática do Mapeamento Periódico}

% VERIFY: Essa seção não deveria vir antes da última (ou substituir ela), pois me parece que ela apresenta as mesmas informações

Considere o cilindro de parede espessa definido no intervalo radial $r \in [R_{\text{int}}, R_{\text{ext}}]$, cuja espessura total de parede é $t = R_{\text{ext}} - R_{\text{int}}$. O cilindro é composto por $N_{\text{cell}}$ células periódicas unitárias (RVEs) \cite{bakhvalov1989}.

A espessura de uma célula periódica unitária (período $Y$) é calculada a partir do número de células e da espessura total do vaso por:

\begin{equation}
  L_c = \frac{R_{\text{ext}} - R_{\text{int}}}{N_{\text{cell}}}
  \label{eq:L_c}
\end{equation}

Sendo $V_1$ e $V_2 = 1 - V_1$ as frações volumétricas dos Materiais 1 (Aço) e 2 (Alumínio), respectivamente, as espessuras individuais de cada subcamada dentro de uma célula são expressas por:

\begin{equation}
  h_1 = V_1 L_c , \quad h_2 = V_2 L_c
  \label{eq:h_camadas}
\end{equation}

A partição espacial radial do cilindro é dividida de forma discreta em um vetor de raios interfaciais $r_k$, pré-calculado e alinhado de forma exata com as coordenadas de nós correspondentes da malha espacial. Para cada célula periódica $i \in \{0, 1, \dots, N_{\text{cel}} - 1\}$, as posições das interfaces são calculadas algoritmicamente por:

\begin{itemize}
  \item Interface interna da célula $i$ (Material 1 para Material 2):
  \begin{equation}
    r_{(2i+1)} = R_{\text{int}} + i L_c + h_1
    \label{eq:interface_interna}
  \end{equation}

  \item Interface de contorno entre a célula $i$ e a célula $i+1$ (Material 2 para o Material 1 da célula subsequente):
  \begin{equation}
    r_{(2i+2)} = R_{\text{int}} + (i + 1) L_c
    \label{eq:interface_externa}
  \end{equation}
\end{itemize}

onde o índice de célula $i$ estende-se até $N_{\text{cel}} - 2$ para a interface de transição de células, uma vez que a última face externa coincide com o limite do raio externo $R_{\text{ext}}$ do cilindro. Essa formulação genérica permite simular de forma robusta e rigorosa frações volumétricas quaisquer das fases.

A função degrau discreta que governa a distribuição espacial do Módulo de Young local $E(r)$ e do Coeficiente de Poisson local $\nu(r)$ ao longo da parede do cilindro baseia-se na identificação de qual intervalo de subcamada indexada $k$ contém a coordenada do ponto de integração Gaussiana $r = x[0]$:

\begin{equation}
  E(r) = \begin{cases} 
    E_1, & \text{se } k \text{ é par (Aço)} \\ 
    E_2, & \text{se } k \text{ é ímpar (Alumínio)} 
  \end{cases}
  \label{eq:mapeamento_E_v2}
\end{equation}

\begin{equation}
  \nu(r) = \begin{cases} 
    \nu_1, & \text{se } k \text{ é par (Aço)} \\ 
    \nu_2, & \text{se } k \text{ é ímpar (Alumínio)} 
  \end{cases}
  \label{eq:mapeamento_nu_v2}
\end{equation}

\subsection{Algoritmo de Implementação via \texttt{UserExpression} no FEniCS}

No ambiente computacional FEniCS, o mapeamento espacial contínuo expresso pelas \eqref{eq:mapeamento_E_v2} e \eqref{eq:mapeamento_nu_v2} é implementado de forma nativa e extremamente eficiente através de subclasses da classe base \texttt{UserExpression} \cite{logg2012automated}. Esta abordagem baseada em pesquisa indexada permite que as propriedades elásticas sejam avaliadas em tempo real durante o processo de montagem da matriz de rigidez global, eliminando gargalos de memória.

A classe customizada desenvolvida para este propósito, denominada \texttt{PeriodicProperty}, é estruturada em Python conforme a rotina real de cálculo implementada no script:

\begin{verbatim}[language=Python, caption={Implementação da classe PeriodicProperty para mapeamento periódico das propriedades elásticas no FEniCS.}]
class PeriodicProperty(UserExpression):
    def __init__(self, val1, val2, interfaces, **kwargs):
        super().__init__(**kwargs)
        self.val1, self.val2 = val1, val2 # Propriedades elásticas das fases
        self.interfaces = interfaces      # Vetor de coordenadas das interfaces radiais

    def eval(self, values, x):
        idx = 0                           # Inicializa o índice da camada como par
        # Varredura sequencial para localizar a camada correspondente
        for r_int in self.interfaces:
            if x[0] <= r_int:
                break
            idx += 1

        # Atribuição condicional baseada na paridade do índice do segmento
        values[0] = self.val1 if idx % 2 == 0 else self.val2

    def value_shape(self):
        return () # Retorna um escalar (campo de propriedades elásticas)
\end{verbatim}

Durante a execução do solver, as instâncias de \texttt{PeriodicProperty} são utilizadas para construir os campos escalares do Módulo de Young e do Coeficiente de Poisson no domínio discreto, escrevendo-se:

\begin{lstlisting}[language=Python, caption={Instanciação dos campos de propriedades no script FEniCS.}]
E_func = PeriodicProperty(E1, E2, interfaces, degree=0)
nu_func = PeriodicProperty(nu1, nu2, interfaces, degree=0)
\end{lstlisting}

Adota-se o parâmetro \texttt{degree=0} (interpolação constante por trecho) para garantir que a propriedade mude de forma abrupta e descontínua exatamente nas coordenadas das interfaces físicas, simulando com fidelidade a transição bimetálica real sem suavizações artificiais indesejadas.

\subsection{Transição Sem Costura entre o Modelo Detalhado e o Modelo Homogeneizado (MHA)}

A grande elegância metodológica desta arquitetura de programação é a facilidade de transição entre o modelo microscópico detalhado (altamente heterogêneo e oscilatório) e o modelo macroscópico equivalente homogeneizado pelo MHA.

No FEniCS, as formas bilinear $a(u,v)$ e linear $L(v)$ são acopladas ao tensor constitutivo de quarta ordem $\mathcal{C}$ \cite{logg2012automated}. Para alternar as análises mecânicas, o operador computacional simplesmente substitui o argumento de rigidez elástica:

1. \textbf{Modelo Detalhado (MEF-D):} O tensor constitutivo local $\mathcal{C}(r)$ é calculado em cada ponto de integração a partir dos campos oscilatórios $E(r)$ e $\nu(r)$ fornecidos pela classe \texttt{PeriodicProperty}, gerando o comportamento real em zigue-zague dos deslocamentos e as oscilações de tensão.
2. \textbf{Modelo Homogeneizado (MHA):} O tensor constitutivo é inicializado diretamente como uma matriz ortotrópica constante $\hat{\mathcal{C}}$, cujos coeficientes efetivos ($\hat{\mathcal{C}}_{11}$, $\hat{\mathcal{C}}_{12}$, $\hat{\mathcal{C}}_{22}$ e $\hat{G}$) são calculados analiticamente através das equações de Backus deduzidas no Capítulo 3 e aplicadas no script.

Esta simplicidade operacional e algorítmica demonstra o elevado nível de maturidade computacional da simulação, reduzindo o tempo de preparação do modelo e viabilizando estudos paramétricos de otimização estrutural de vasos de pressão multicamadas.

\section{Propriedades dos Materiais e Cenários de Carga}\label{sec:propriedades_cenarios}

Para a execução das simulações numéricas via Método dos Elementos Finitos (MEF) e a consequente avaliação da precisão do Método de Homogeneização Assintótica (MHA), torna-se indispensável definir com precisão os parâmetros físicos e mecânicos que caracterizam as fases constituintes do cilindro bimetálico. Esta seção detalha as propriedades elásticas, as densidades volumétricas, os limites de escoamento adotados para o aço e o alumínio, bem como justifica a escolha do cenário de carregamento termomecânico imposto ao modelo.

\subsection{Propriedades Mecânicas e Físicas das Fases Constituintes}

O cilindro periódico analisado é composto pelo par metálico bimetálico formado por aço carbono (ASTM A516 Gr. 70) e liga de alumínio 6061-T6, ambos tratados individualmente como meios homogêneos e isotrópicos em nível de subcamada microestrutural. As propriedades mecânicas de entrada (módulo de Young $E$ e coeficiente de Poisson $\nu$), associadas às propriedades físicas de massa específica ($\rho$) e resistência mecânica elástica limite ($\sigma_y$), são sumarizadas na Tabela~\ref{tab:propriedades_fases}:

\begin{table}[htbp]
  \centering
  \caption{Propriedades elásticas, físicas e limites de resistência das fases constituintes do cilindro bimetálico conforme implementação no script de projeto.}
  \label{tab:propriedades_fases}
  \begin{tabular}{|l|c|c|c|}
    \hline
    \textbf{Propriedade} & \textbf{Símbolo} & \textbf{Fase 1: Aço} & \textbf{Fase 2: Alumínio} \\ \hline
    Módulo de Elasticidade & $E$ & $200 \, \text{GPa}$ ($2,0 \times 10^{11} \, \text{Pa}$) & $70 \, \text{GPa}$ ($7,0 \times 10^{10} \, \text{Pa}$) \\ \hline
    Coeficiente de Poisson & $\nu$ & $0,30$ & $0,33$ \\ \hline
    Densidade Volumétrica & $\rho$ & $7800 \, \text{kg/m}^3$ & $2700 \, \text{kg/m}^3$ \\ \hline
    Limite de Escoamento & $\sigma_y$ & $260,0 \, \text{MPa}$ ($2,6 \times 10^8 \, \text{Pa}$) & $270,0 \, \text{MPa}$ ($2,7 \times 10^8 \, \text{Pa}$) \\ \hline
  \end{tabular}
\end{table}

A escolha desses materiais reflete um cenário clássico de alto contraste de rigidez elástica (com a razão de módulos $E_{\text{aço}}/E_{\text{alu}} \approx 2,86$), o qual representa o cenário ideal e mais rigoroso para testar a capacidade do MHA de suavizar as oscilações de alta frequência espacial sem perda de consistência energética.

A massa linear $M_{\text{lin}}$ (expressa em kg por metro de comprimento axial do cilindro) é calculada diretamente a partir da área da seção transversal da parede do vaso $A_{\text{wall}} = \pi (R_{\text{ext}}^2 - R_{\text{int}}^2)$ e da densidade média do material. Para os raios de projeto ($R_{\text{int}} = 9,0 \, \text{mm}$ e $R_{\text{ext}} = 11,0 \, \text{mm}$), a área do resulta em $A_{\text{wall}} \approx 125,66 \, \text{mm}^2$. A densidade equivalente macroscópica ($\hat{\rho}$) do compósito com fração volumétrica simétrica ($V_1 = V_2 = 0,5$) é regida estritamente pela regra aditiva simples:

\begin{equation}
  \hat{\rho} = V_1 \rho_{\text{aço}} + V_2 \rho_{\text{alu}} = 0,5 (7800) + 0,5 (2700) = 5250 \, \text{kg/m}^3
  \label{eq:densidade_efetiva}
\end{equation}

% TODO: Adicionar uma fonte que justifique o uso da regra aditiva simples, mostrando suas limitações

Resultando em uma massa linear de $0,6597 \, \text{kg/m}$ para o cilindro bimetálico, o que representa um alívio de peso de exatamente $32,69\%$ em comparação a uma estrutura fabricada exclusivamente em aço maciço ($0,9802 \, \text{kg/m}$).%, confirmando a elevada eficiência de peso da hibridização de materiais.

% REMOVE: Última afirmação removida por receio quanto à viabilidade de manufatura desse compósito.

\subsection{Cenários de Carregamento e Pressão de Projeto}

Para submeter o modelo numérico axissimétrico e a formulação analítica do MHA a condições operacionais realistas e severas, estabeleceu-se uma \textbf{pressão interna de projeto estável de $30 \, \text{MPa}$} ($3,0 \times 10^7 \, \text{Pa}$) aplicada de forma contínua na fronteira radial interna do cilindro ($r = R_{\text{int}}$).

Esta escolha de carregamento situa-se no patamar de pressões de serviço reportado para vasos de pressão Tipo III de armazenamento de hidrogênio, para os quais a literatura indica pressão de projeto inferior a $35 \, \text{MPa}$ \cite{Zhang2022}.

A imposição deste patamar de solicitação induz um estado de tensões internas elevadas na parede do cilindro de parede grossa ($R_{\text{int}}/t = 4,5$), permitindo-nos avaliar:
1. \textbf{O escoamento inicial das fases:} Onde o aço ASTM A516 Gr. 70 (com escoamento de $260,0 \, \text{MPa}$) e o alumínio (escoamento de $270,0 \, \text{MPa}$) são testados sob as tensões equivalentes locais de Von Mises calculadas pelo MEF;
2. \textbf{A precisão microscópica do MHA:} Provando se a solução analítica homogeneizada da modelagem assintótica consegue prever as tensões médias e as deformações globais em regime elástico linear mesmo sob pressões internas severas de $300 \, \text{bar}$ ($30 \, \text{MPa}$).

As condições de contorno de Neumann de pressão interna foram aplicadas no contorno radial interno $\Gamma_{\text{int}}$ por meio da forma linear de trabalho virtual descrita na Seção 3.4, garantindo a perfeita integração espacial ponderada pelo raio $r$ exigida pelo FEniCS \cite{logg2012automated}.

% TODO: Tornar número da seção 3.4 dinâmico